\clearpage
\section{Introduction}

\textbf{Aerial Inspection in Confined Spaces.}
The deployment of Unmanned Aerial Vehicles (UAVs) has significantly advanced the field of industrial inspection. In expansive, outdoor environments, drones are now routinely utilized to inspect power lines, bridge infrastructure, and wind turbines, drastically reducing the risk to human operators while increasing efficiency. However, the transition of autonomous aerial robotics into indoor, confined spaces remains a significant challenge. Industries rely heavily on complex internal structures---such as large-scale storage vats, subterranean networks, and industrial piping---that require frequent inspection. Yet, the adoption of autonomous UAVs in these spaces severely lags behind outdoor applications due to a severe set of environmental constraints \cite{ozaslan2017inspection}.

\textbf{Navigation Challenges in GPS-Denied Environments.}
Operating a UAV within confined industrial spaces, such as a stainless steel storage tank, fundamentally alters the requirements for flight. These environments are routinely GPS-denied and, due to heavy ferromagnetism or physical shielding, completely signal-denied. This severs the link to human pilots, mandating full onboard autonomy. However, achieving autonomy in these conditions exposes the limitations of standard sensor suites. Featureless or highly reflective surfaces confound optical cameras, while traditional 3D LiDAR systems are often too heavy for the payload capacity of a drone small enough to enter the space \cite{khattak2020robust}. Furthermore, even if the sensors could accurately map the environment, poses a substantial computational burden---a challenge that even ground-based autonomous vehicles struggle to solve reliably in much simpler 2D environments \cite{macariorojas2021}.

\textbf{Collision Resilience as a Navigation Modality.}
Given the strict payload limits on onboard companion computers and the unreliability of traditional sensors in hostile environments, complex, real-time collision avoidance becomes computationally impractical. In tight spaces, collisions are not a symptom of system failure; they are an inevitability. Therefore, the paradigm of autonomous confined-space flight must shift from attempting to avoid obstacles at all costs to surviving them and, ultimately, exploiting them.

Surviving these inevitable impacts requires a mechanical design capable of significantly lowering the kinetic penalty of a collision. This could perhaps be achieved by a structural design that, to a certain extent, deflects the collision forces away from the drone's main frame. While commercial solutions addressing this exist---such as the Flyability Elios \cite{briod2014collision}, which utilizes a freely rotating cage to roll along surfaces---empirical flight data detailing their performance during impact events remain sparse in the open literature---a gap that must be investigated before physical impacts can be reliably exploited for navigation.

\textbf{Research Objectives.}
If physical collisions in signal-denied environments are inevitable, and if a UAV can be designed to reliably survive them, a new avenue for spatial perception emerges. Rather than relying on heavy optical or laser-based sensors, a drone could theoretically utilize the very impacts it sustains to understand its surroundings. By interpreting the raw inertial data generated at the moment of collision, an autonomous system could extract foundational geometric information about its environment. If a drone operating without an external localization system can estimate the geometric angle of an impact using only its IMU, it gains a form of ``tactile odometry.'' Successive collisions could allow the system to map the boundaries of a confined space purely by bouncing off them. This transition---from a drone that is merely blind and resilient to one that is blind but tactile---forms the core of this investigation.

\subsection{Research Questions}

To address the need to deflect impact forces and redirect linear momentum, and to investigate the potential for collision exploitation, this thesis details the physical and software development of a custom, autonomous quadcopter equipped with an onboard companion computer and a freely rotating protective cage. By mechanically decoupling the protective outer shell from the drone's main frame, this design introduces a mechanism to redirect linear momentum into cage rotation, mitigating the impact penalty transferred to the airframe.

Using a ground truth motion capture (MoCap) environment to validate the system, this research seeks to answer the following questions:

\begin{enumerate}
    \item How does the physical resilience and energy deflection capability of a freely rotating protective cage compare to a fixed cage during direct obstacle collisions?
    \item To what extent can the raw inertial data generated during these deliberate physical collisions be utilized to accurately estimate the angle of impact against an obstacle, serving as a fundamental step toward tactile spatial awareness in signal-denied environments?
\end{enumerate}
